\chapter{Terminologia específica}
\label{cap:glossario}

A \cref{tab:glossario} reúne, por ordem alfabética, os termos técnicos
usados neste relatório. Quando o termo é habitualmente usado em inglês,
indica-se o equivalente entre parênteses.

{\small
\begin{longtable}{@{}L{4.4cm}L{10.0cm}@{}}
  \caption{Terminologia específica}
  \label{tab:glossario}\\
  \toprule
  \textbf{Termo} & \textbf{Definição} \\
  \midrule
  \endfirsthead
  \toprule
  \textbf{Termo} & \textbf{Definição} \\
  \midrule
  \endhead
  \bottomrule
  \endfoot
  ADAS
    & Sistemas avançados de assistência à condução (\emph{Advanced Driver
      Assistance Systems}), como a travagem automática de emergência ou a
      manutenção na faixa. \\
  AEB
    & Travagem automática de emergência (\emph{Autonomous Emergency
      Braking}): função que trava o veículo quando deteta uma colisão
      iminente. \\
  Agenda
    & Em Drools, conjunto das regras ativadas, cujas condições estão
      satisfeitas, que aguardam execução. A ordem de execução é
      determinada pela estratégia de resolução de conflitos. \\
  Aquisição de conhecimento
    & Processo de obter, junto de peritos e de outras fontes, o
      conhecimento necessário para construir a base de conhecimento de um
      sistema pericial. \\
  Base de conhecimento
    & Componente de um sistema pericial que contém o conhecimento do
      domínio sob a forma de factos, regras, prioridades e exceções. \\
  Confiança da deteção
    & Medida do grau de certeza com que um sensor, ou a fusão de vários,
      identifica um objeto ou uma condição. \\
  Desempate
    & Critério aplicado quando duas ou mais ações candidatas têm a mesma
      prioridade. Neste projeto, escolhe-se a primeira de uma lista
      ordenada de ações (\cref{sec:alteracoes-bc}). \\
  Distância de paragem
    & Distância percorrida pelo veículo desde a deteção de um perigo até à
      imobilização. Inclui a distância percorrida durante o tempo de
      reação e a distância de travagem. \\
  Drools
    & Motor de regras de produção de código aberto para Java, usado na
      implementação do sistema. \\
  Encadeamento progressivo
    & Estratégia de inferência que parte dos factos conhecidos e aplica as
      regras cujas condições estão satisfeitas para obter novas conclusões
      (\emph{forward chaining}). \\
  Encadeamento regressivo
    & Estratégia de inferência que parte de uma conclusão a provar e
      procura os factos e regras que a sustentam (\emph{backward
      chaining}). \\
  Explicação \emph{How}
    & Explicação que mostra a cadeia de factos e regras que conduziu a uma
      conclusão. \\
  Explicação \emph{Why}
    & Explicação que justifica por que motivo uma ação foi escolhida,
      indicando a regra aplicada e as condições satisfeitas. \\
  Explicação \emph{Why not}
    & Explicação que justifica por que motivo uma ação alternativa não foi
      escolhida, indicando as condições não satisfeitas ou a ação de maior
      prioridade que prevaleceu. \\
  Facto
    & Unidade de informação sobre a situação atual, mantida na memória de
      trabalho; por exemplo, \facto{obstaculo\_confirmado}. \\
  Falso positivo
    & Deteção de um objeto ou condição que não existe na realidade,
      resultante, por exemplo, de uma leitura atípica de um sensor. \\
  Fusão sensorial
    & Combinação de dados de vários sensores para obter uma representação
      do ambiente mais fiável do que a de cada sensor isolado
      (\emph{sensor fusion}). \\
  GNSS
    & Sistema global de navegação por satélite (\emph{Global Navigation
      Satellite System}), de que o GPS é um exemplo. \\
  IMU
    & Unidade de medição inercial (\emph{Inertial Measurement Unit}):
      sensor que mede acelerações e velocidades angulares do veículo. \\
  Intervalo temporal
    & Tempo que o veículo próprio demora a percorrer a distância que o
      separa do veículo precedente, à velocidade atual (\emph{time
      headway}). \\
  LiDAR
    & Sensor que mede distâncias através de impulsos laser e produz uma
      nuvem de pontos do ambiente (\emph{Light Detection and
      Ranging}). \\
  Limiar
    & Valor de referência a partir do qual uma condição de uma regra é
      considerada satisfeita (\emph{threshold}). \\
  Memória de trabalho
    & Componente de um sistema pericial que contém os factos relativos à
      situação atual, avaliados pelo motor de inferência (\emph{working
      memory}). \\
  Módulo de explicação
    & Componente de um sistema pericial que justifica as conclusões
      produzidas (\emph{explanation facility}). \\
  Motor de inferência
    & Componente de um sistema pericial que avalia as regras da base de
      conhecimento sobre os factos da memória de trabalho e produz
      conclusões (\emph{inference engine}). \\
  Níveis de automação
    & Classificação da norma SAE J3016, de 0 (sem automação) a 5
      (automação integral), segundo a parte da tarefa de condução assumida
      pelo sistema. \\
  Paragem imediata
    & Ação de segurança de prioridade máxima, aplicada em caso de falha do
      veículo ou de dados contraditórios numa situação de nível alto ou
      crítico (\emph{STOP}). \\
  Perito
    & Pessoa com conhecimento especializado do domínio, consultada durante
      a aquisição de conhecimento. \\
  Phreak
    & Algoritmo de correspondência de padrões usado nas versões atuais do
      Drools, derivado do algoritmo Rete. \\
  Ponto de não-retorno
    & Instante a partir do qual a travagem já não permite evitar a
      colisão. \\
  Prioridade
    & Nível de importância associado a uma regra ou ação: baixo, médio,
      alto ou crítico. \\
  Radar
    & Sensor que usa ondas de rádio para medir a distância e a velocidade
      relativa de objetos. \\
  Regra de produção
    & Regra da forma SE condições ENTÃO ações, usada para representar
      conhecimento num sistema pericial (\emph{production rule}). \\
  Resolução de conflitos
    & Processo de escolha entre várias regras ou ações aplicáveis em
      simultâneo. \\
  Rete
    & Algoritmo eficiente de correspondência entre muitos padrões e muitos
      factos, proposto por Forgy~\cite{forgy1982} e na base de motores de
      regras como o Drools. \\
  Saliência
    & Atributo de uma regra em Drools que define a sua prioridade relativa
      na agenda (\emph{salience}). \\
  Sensor exterocetivo
    & Sensor que observa o ambiente em redor do veículo, como a câmara, o
      radar ou o LiDAR. \\
  Sensor propriocetivo
    & Sensor que mede o estado do próprio veículo, como a velocidade das
      rodas, as acelerações ou o ângulo da direção. \\
  Sessão com estado
    & Sessão do motor de regras que mantém os factos na memória de
      trabalho entre ciclos de inferência (\emph{stateful session}). \\
  Sistema pericial
    & Sistema que representa conhecimento especializado de forma explícita
      e o utiliza, através de um motor de inferência, para chegar a
      conclusões justificáveis (\emph{expert system}). \\
  Tempo até à colisão (TTC)
    & Tempo que falta para uma colisão se as velocidades se mantiverem:
      distância ao objeto dividida pela velocidade de aproximação
      (\emph{time to collision}). \\
  TTC
    & Sigla de \emph{time to collision}; ver tempo até à colisão. \\
  Veículo precedente
    & Veículo que circula imediatamente à frente do veículo próprio, na
      mesma faixa (\emph{lead vehicle}). \\
  Veículo próprio
    & Veículo controlado pelo sistema (\emph{ego vehicle}). \\
  Velocidade relativa
    & Diferença entre a velocidade do veículo próprio e a de outro objeto,
      medida na mesma direção. \\
\end{longtable}
}
